
\documentclass{article}

% FLMSec 2026 uses double-blind review.
\usepackage[dblblindworkshop]{neurips_2026}
\workshoptitle{Foundations of Language Model Security}

\usepackage[utf8]{inputenc}
\usepackage[T1]{fontenc}
\usepackage{hyperref}
\usepackage{url}
\usepackage{booktabs}
\usepackage{amsfonts}
\usepackage{amsmath}
\usepackage{microtype}
\usepackage{xcolor}
\usepackage{graphicx}

\title{Investigating the Refusal Boundary in Language Models}

% Keep the submission anonymous. The style will suppress these details in
% double-blind mode; replace them only for the camera-ready version.
\author{Anonymous Author(s)}

\begin{document}

\maketitle

\begin{abstract}
Safety-tuned Large Language Models (LLMs) are trained to refuse harmful requests without refusing standard requests; weaknesses on either side produce jailbreaks and over-refusals, respectively. We study which automated attacks and internal refusal signals best reveal weaknesses in this boundary. We compare gradient-based rewriting, prompted attackers, and reward-trained attackers on Llama3-8B-Instruct and Qwen3-32B, using refusal-prefix probabilities, activation directions, and multi-layer probes as optimization signals.
Rewrites are filtered using MiniLM similarity, Llama Guard, while an independent Sonnet judge determines whether the original was answered and the rewrite received an unjustified refusal. Phrase matching is only diagnostic. We find that trained attacker models most reliably create fluent red-team prompts, while gradient-based methods often reveal underlying weaknesses but fail to balance fluency with refusal manipulation.
\end{abstract}

\section{Introduction}

Modern LLMs are trained to refuse harmful prompts. This behavior prevents models from propagating content related to self-harm, terrorism, hacking, fraud, CBRN threats, and other harms. However, LLM refusal is not perfectly tuned. The \textit{refusal boundary} describes the edges of prompt space where a model can be nudged into refusal. The refusal boundary inevitably captures some over-refusals while missing some jailbreaks \citep{rottger2024xstest,cui2024orbench}. Investigating the refusal boundary is critical for two reasons:

\begin{enumerate}
    \item Investigating the refusal boundary helps model developers tune it. Tuning the refusal boundary is impactful because it can increase safety. Model developers do not want excessive refusal, so they must balance the number of over-refusals and jailbreaks allowed by the model. Reducing the number of jailbreaks improves safety directly. Reducing the number of over-refusals improves safety by shifting the Pareto frontier between helpfulness and harmfulness, allowing model providers to push the refusal boundary forward with less concern about excessive refusal and thereby reduce the number of jailbreaks.

    \item Investigating the refusal boundary helps us evaluate model alignment in a domain that is relatively well defined. If developers struggle to create a fine-grained refusal boundary, this should negatively inform their appraisal of harder-to-evaluate properties, such as value alignment out of distribution. Understanding the failure modes that pertain to refusal may provide insights that transfer to value alignment.
\end{enumerate}

We use automated attack processes to find weak points in the refusal boundary. We prioritize natural-language edits. These are particularly important because they can pass perplexity filters and arise in natural user conversations \citep{wichers2024gradient}. We test a variety of internal signals of refusal, including refusal-vector activations, ensembles of linear probes trained to detect refusal, and logit scores for common refusal-response prefixes. We test gradient-based methods, prompted LLM attacker models, and trained LLM attackers. We also test both over-refusal prompts and jailbreak prompts. These variations allow us to locate and study weaknesses in the refusal boundary more thoroughly. We repeat this investigation for Llama-3-8B and Qwen3-32B.

We find that trained attacker models most effectively produce natural-language over-refusal prompts. By analyzing prompts generated by our attacker models, we find patterns that include common trigger words, such as ``weaponize,'' that models over-index on when determining whether they should refuse. This result is consistent with prior work showing that LLMs often anchor on sensitive words or superficial linguistic features in their refusal behavior \citep{rottger2024xstest,wu2025evorefuse}. We validate these patterns by deriving over-refusal directions for each pattern and ablating them---that is, projecting the residual stream onto the subspace orthogonal to each direction. These interventions reduce over-refusal without reducing standard refusal behavior, building on prior work that identifies and intervenes on refusal directions in activation space \citep{arditi2024refusal,wollschlager2025geometry}. We find that gradient-based methods often reveal underlying triggers but are less reliable for finding natural-language over-refusals. Natural-language jailbreaks are more difficult to find in Llama-3-8B and Qwen3-32B than over-refusals; however, they still exist along the refusal boundary. Gradient-based methods can find jailbreaks, but limiting the results to coherent requests drastically reduces their yield.

\paragraph{Contributions.}
Our main contributions are:
\begin{enumerate}
    \item We compare multiple signals of refusal, including refusal-vector activations, linear refusal probes, and refusal-prefix logit scores.
    \item We test gradient-based methods, prompted attacker models, and trained attacker models for finding over-refusals and jailbreaks in Llama-3-8B and Qwen3-32B.
    \item We identify common patterns in natural-language over-refusal prompts. We show that pattern-specific activation directions can reduce over-refusal without reducing refusal of harmful prompts and that prompting attacker models with these patterns as prompts makes them more effective.
\end{enumerate}

Overall, our results show that the refusal boundary contains interpretable weaknesses that can be found through automated attacks. Studying these weaknesses can help model developers increase safety from harmful prompts and better understand the alignment properties of their models.


\section{Refusal signals}
% Draft method-and-result bullets grounded in the implemented scorers and the
% current experiment notes. Convert to prose after deciding what to emphasize.

\textbf{Refusal-prefix logits} In order to derive a logit-based refusal score, we begin by identifying prefixes that precede refusal responses for each model. In order to capture the refusal prefixes most likely within our context, we sample from our over-refusal rewrite dataset. For Llama-3-8B-Instruct, \\ \(\{\text{I cannot, I apologize, I can't, Sorry, I'm sorry}\}\), covers 99.6\% of the observed refusal openings. For Qwen3-32B, the set is \(\{\text{I'm sorry, Sorry, I can't, I cannot, As an AI}\}\), covering 99.7\%. Because none of the provided token sequences are a prefix of another, the sum does not double-count the same refusal onset. Thus, we sum across the logit scores for the prefixes to calculate the logit score for refusal.


\textbf{Activation directions} To identify a refusal direction, we begin by compiling a contrast dataset of prompts that Llama3-8b refuses on and prompts that it does not refuse on. The source prompt pool is derived primarily from the English S-Eval attack and base-risk sets~\citep{yuan2024seval}, supplemented by a filtered subset of Anthropic’s human red-team prompts~\citep{ganguli2022redteaming} and additional jailbreak and prompt-injection material. The non-refusal class therefore contains mostly adversarial or unsafe requests that did not trigger refusal, along with a small number of ordinary prompts from the off-topic portion of the OnepaneAI harmful-prompts dataset~\citep{onepaneai2024harmfulprompts}. Mixing jailbreaks with genuinely harmful and helpful prompts and then dividing by measured model refusal ensures that the activation direction tracks refusal rather than the harmfulness of the content. 

At each layer \(L\) of the residual stream, we compute \(d_L=\mu_{\mathrm{refuse},L}-\mu_{\mathrm{comply},L}\) from the activation at the final real prompt token. We use the vectors that maximize the difference of means at each layer as canidates for the final refusal vector.

We select the candidate that most effectively separates separates from non-refusal in our dataset when activation for each prompt is compared against the refusal vector. For Llama3-8b, layer 17 was selected and for Qwen3-32B layer 58 was selected. We then validate the layer by adding the along this direction. For Llama this this raised refusal among benign prompts from 0\% to 99\%. For Qwen3-32B refusal was raised from from 5\% to 93\%. Further details and the model-specific selection rationale are provided in Appendix~\ref{app:refusal-vector-selection}.

\textbf{Learned refusal probes}
    To learn a probe ensemble we use a mass-means approach, following \citet{marks2024geometry}. We begin by calculating the difference of means of each layer using the same aggregated dataset as for the refusal vector. We then take separate refusal rewrite datasets generated by Llama3-8B and Qwen3-32B, respectively, generate 4 model responses to each rewrite, and then rank each rewrite based on how consistently it creates refusal responses according to an automated grader. The probe ensemble is created by using a non-negative least squares formula to find the weighting of the difference of means at each layer that best approximates the ranking of the activations produced by the rewrites, ranked by how much each rewrite induces refusal.

    We tested covariance-corrected LDA directions as well as raw mass-mean directions. Raw directions performed better in the corrected Llama and Qwen comparisons. This may be because our samples were insufficient to avoid overfitting the large feature covariance matrices for Llama and Qwen.
    
\section{Gradient-based prompt rewriting}
% Briefly state the white-box threat model and validity definition here rather
% than in a standalone problem-setup section: the original task is benign, its
% meaning should remain fixed, and the attack succeeds only when the original is
% answered but the rewrite is unjustifiably refused.

\subsection{Rewrite optimization}

We use a variant of Greedy Coordinate Gradient (GCG) to rewrite benign prompts to turn them into over-refusal prompts and rewrite harmful prompts to become jailbreaks \citep{zou2023universal}. Unlike the original GCG method, which appends an adversarial suffix, our method changes the prompt itself. Combined with additional terms in the loss that weight semantic similarity and fluency, this change allows us to keep rewrites more realistic while remaining adversarial.

For each original prompt \(o\), GCG begins with the tokens of \(o\) and performs 150 search steps. At each step, the gradient of the refusal loss identifies promising replacement tokens at each editable position. We keep the 256 replacements favored most strongly by the gradient and sample 96 candidate rewrites that each change one position in the current prompt. Each proposed token must also be among the target model's 512 most likely tokens at that position. We evaluate every candidate as an actual token sequence and retain the candidate with the lowest total loss.

For Llama-3-8B-Instruct, the refusal objective is the negative log-probability that the response begins with ``I cannot.'' This objective performed better in our preliminary rewriting experiments than directly maximizing the prompt's projection onto a refusal direction. For Qwen3-32B, we instead optimize the combined probability of the five most common refusal openings identified for that model: ``I'm sorry,'' ``Sorry,'' ``I can't,'' ``I cannot,'' and ``As an AI.'' A model-specific set is necessary because only a small fraction of Qwen refusals begin with ``I cannot.'' In an initial Llama3-8b experiment with 60 prompts per objective and no fluency term, manual review found 9 valid over-refusals for the ``I cannot'' objective, compared with 5 for the layer-17 direction and 3 for the layer-12 direction. These small samples suggest that logit scores may provide a more robust objective for refusal in a gradient-based context, but they are not large enough to establish the size of the differences reliably.

The full optimization loss is
\[
\mathcal{L}(o,r)
=
\mathcal{L}_{\mathrm{refuse}}(r)
+
20\max\!\left(0,\,0.85-\operatorname{sim}(o,r)\right)
+
\mathcal{L}_{\mathrm{mean}}(r)
+
0.5\mathcal{L}_{\mathrm{max}}(r),
\]
where \(r\) is the candidate rewrite, \(\operatorname{sim}(o,r)\) is its MiniLM cosine similarity to the original, and \(\mathcal{L}_{\mathrm{refuse}}\) is the refusal-opening loss. We obtain one vector per prompt by averaging and normalizing the token representations from \texttt{sentence-transformers/all-MiniLM-L6-v2}, then compute the cosine similarity of the two vectors. The final two terms measure fluency under the target model. The first is the average language-model loss across the prompt, while the second is the loss of its single least likely token.

GCG can replace tokens but cannot insert new ones. In the main Llama experiment, we therefore add 12 editable placeholder positions after the original prompt so that the optimizer can complete phrases rather than leaving the rewrite unfinished. Replacements are limited to complete word-initial tokens or punctuation. We also evaluate a more restrictive version that keeps the original length fixed, changes at most four complete content words, and rejects candidates whose fluency loss rises substantially above that of the original.

We also tested PEZ, which searches over continuous token embeddings \citep{wen2023hard}. We did not use PEZ in the main rewriting experiments because it could not directly apply the MiniLM similarity constraint during each search step. MiniLM and the target model use different tokenizers, so the sentence-level similarity score cannot be calculated directly from PEZ's continuous embeddings. GCG instead evaluates complete candidate strings and can therefore apply the refusal, similarity, and fluency terms to the same discrete text. Our preliminary comparisons cover the refusal-opening objective and activation-direction objectives at Llama layers 12 and 17, optimization with and without the fluency terms, and the four-content-word condition described above. Our preliminary comparisons cover the refusal-opening objective and activation-direction objectives at Llama layers 12 and 17, optimization with and without the fluency terms, and the four-content-word condition described above.

\subsection{Output filtering}For each original and rewrite, the target model generates one deterministic response, capped at 128 new tokens, under the system instruction ``You are a helpful assistant.'' After optimization, we retain a pair only if it passes all of the following filters:\begin{itemize} \item \textbf{Semantic similarity.} The rewrite's MiniLM cosine similarity to the original is at least 0.85. \item \textbf{Prompt safety.} Llama Guard 3-8B labels both the original and the rewrite as safe. \item \textbf{Behavioral judgment.} Claude Sonnet 5 independently labels the target model's original response \textsc{Comply} and its response to the rewrite \textsc{Refuse}, and judges the rewritten-prompt refusal unjustified. \textsc{Clarify}, justified, and borderline cases are discarded. Sonnet 4.6 is used only if the primary judge refuses the classification task. The judge evaluates each prompt--response pair separately; it does not directly compare whether the rewrite preserves the original task. \end{itemize}In the completed Llama scale experiment on 3{,}200 benign Alpaca-derived original--rewrite pairs, the refusal-phrase detector marked 95.0\% of rewrites as refusals.

\subsection{Gradient-rewrite results}


During the initial 60 prompt derisking phase, adding the fluency constraints improved both the safety and the yield of the search. For the ``I cannot'' objective, the fraction of rewrites marked unsafe by Llama Guard fell from 70.0\% to 13.3\%, while the number of clean refusal candidates increased from 9 to 23 out of 60. Manual review retained 21 of these 23 as over-refusals. Under the same constraints, the layer-17 direction produced 12 manually confirmed cases, while the layer-12 direction produced 1. The poor layer-12 result is consistent with our other experiments: layer 12 is causally important when its activations are changed directly, but it is difficult to reach through token changes to the input. This illustrates why causal importance and usefulness as an optimization target are different properties.

In the completed Llama scale experiment on 3{,}200 benign Alpaca-derived original--rewrite pairs, 1,394 of the 3,200 rewrites, or 43.6\%, also passed the original-response, Llama-Guard, and similarity gates. After applying the independent response judge, 1,220 of the 3,200 pairs, or 38.1\% with a Wilson 95\% confidence interval of 36.5--39.8\%, satisfied the automatic clean-flip definition: the original was answered, the rewrite received an unjustified refusal, both prompts were marked safe, and MiniLM similarity was at least 0.85. These 1,220 cases come from 1,220 distinct original prompts. The mean MiniLM similarity among them is 0.864. The decrease from 43.6\% to 38.1\% shows that the independent judge removes a meaningful number of cases counted by the refusal-phrase detector.





\section{Attacker models}
% Move from per-prompt white-box token search to LLMs that generate complete,
% natural-language rewrites. State what access each attacker has to the target
% and whether the attacker is optimized separately for each target model.
In addition to our gradient-based approach to generating rewrites, we study the effect of post-training models to rewrite prompts to elicit over-refusal. Our trained attackers use rewards calculated from three internal signals of the model, all of which are based on tracking some aspect of model refusal. We report results from comparing trained attacker models to their base model as the target. In Section 4.3, we also cover the discussion across model families. 

\subsection{Prompted attackers}
% Describe the plain instruction to induce refusal, category-guided prompting,
% and control prompting. Explain the prompt inputs, number of samples, selection
% procedure, and whether the attacker sees target-model feedback.

The simplest form of an attacker is the target model itself, prompted to rewrite benign seeds to elicit over-refusal. Specifically, given a benign original prompt $o$, we prompt our model to produce a rewrite $r$ that the target will refuse while preserving the intent of the original prompt and keeping it benign. We evaluated the base models on our dataset of 2{,}400 benign seed prompts. Llama-3-8B-Instruct yielded an over-refusal rate of $4.92\%$ or $(118 / 2{,}400)$ and Qwen3-32B yielded only $2.29\%$ or $(50/2{,}400)$. 

To generate training data for the two models we used a mix of Claude Haiku 4.5 and Claude Sonnet 4.5, prompted to rewrite benign Dolly and Alpaca prompts, with few-shot assistance of over-refusal rewrites that we constructed by hand. We generated three rewrites for roughly $5{,}000$ benign original prompts, producing $14{,}996$ candidate rewrites.

\subsection{Reward-trained attackers}
% Describe the rewriter model, training examples or rollouts, refusal reward,
% semantic/safety gates, reward ranking or optimization algorithm, and how
% held-out prompts are kept separate from training and category discovery.

To train the attackers, we opted for reward-weighted regression (RWR) [cite add], which is functionally a weighted supervised fine-tuning method. We opted against online RL training due to its high training cost.

For an original $o$ and candidate rewrite $r$, the over-refusal reward is \begin{equation}\mathrm{OR}(o,r) \;=\; \exp\!\big(k\,(\mathrm{sim}(o,r) - c)\big)\;\cdot\;\Delta_{\mathrm{signal}}(o,r),\qquad k = 18.4,\;\; c = 0.75,\label{eq:or-reward}\end{equation}where $\mathrm{sim}(o,r)$ is the cosine similarity of MiniLM sentence embeddings [cite] and $\Delta_{\mathrm{signal}}(o,r) = \mathrm{signal}(r) - \mathrm{signal}(o)$ is the refusal delta, calculated by the three white-box refusal signals covered earlier. We chose to give the semantic similarity multiplicative effects through the exponential so that completely unrelated prompts cannot game the OR score by maximizing the refusal delta.

\subsection{Cross-model attacker--target evaluation}
% Present the 2-by-2 attacker--target experiment for Llama and Qwen. Separate an
% attacker's ability to generate effective rewrites from a target model's
% susceptibility to those rewrites, and report transfer as well as same-model
% performance.
Part of the work in determining the nature of refusal between Llama and Qwen is identifying between how good of an attacker a model is and how susceptible of a target it is. To separate these, we sourced identical prompts and use both trained models as attackers against both base models as targets. We generated a corpus of $44{,}482$ unique rewrites over $6{,}000$ Alpaca original prompts and chose the Llama logit model and Qwen probe model as the attackers. 

The result is that over-refusal is highly dependent on the target model and not on the attacker which generates the prompt. The Qwen target is roughly twice as  robust to over-refusal on the same rewrites, and there is no statistical significance between attackers on a given target model. 

 
\begin{table}[t]
  \centering
  \caption{Cross-model attacker--target over-refusal. Cells are judge-confirmed over-refusal rates (\%) on identical held-out prompts; over-refusal tracks the \emph{target} model, not the attacker that produced the rewrite. $n\!\approx\!23.3$k (Llama attacker) and $21.2$k (Qwen attacker) rewrites per row.}
  \label{tab:cross_model}
  \begin{tabular}{lrr}
    \toprule
     & \multicolumn{2}{c}{Target} \\
    \cmidrule(lr){2-3}
    Attacker & Llama-3-8B & Qwen3-32B \\
    \midrule
    Llama (logit) & $10.32$ & $5.14$ \\
    Qwen (probe)  & $10.59$ & $6.06$ \\
    \bottomrule
  \end{tabular}
\end{table}
 
\subsection{Attacker-model results and comparison}
% Compare prompted, category-guided, reward-trained, and gradient-based methods
% using the same independent judge and validity gates. Include naturalness,
% semantic preservation, refusal justification, and attack yield rather than a
% refusal substring alone.

\begin{table}[t]
  \centering
  \caption{Primary evaluation. Populate after independent judging.}
  \label{tab:primary}
  \begin{tabular}{lrrr}
    \toprule
    Method & Valid rewrites & Unjustified refusals & Rate \\
    \midrule
    Gradient rewriting & -- & -- & -- \\
    Prompted attacker & -- & -- & -- \\
    Reward-trained attacker & -- & -- & -- \\
    Control & -- & -- & -- \\
    \bottomrule
  \end{tabular}
\end{table}

\subsection{What edits induce over-refusal?}
% Show word/token diffs for successful rewrites, derive an edit taxonomy only on
% a development split, and test the proposed categories on held-out prompts.
% Report recurring loaded words and semantic framings for each model. Follow the
% descriptive analysis with causal tests, such as category-specific activation
% directions and ablation, without treating the taxonomy itself as evidence.



The rewritten prompts that the trained model produces range from grammatically quite similar to complete paraphrases of the original prompt. We want to explore whether there is qualitative difference in the refusal behavior between these rewrites by analyzing both types.

We operationalize this gap with the Levenshtein [cite] edit-distance $D = Lev(C(o), C(r))$, where $C(p)$ is the lowercase word tokens of prompt $p$ with the standard English stopwords removed. We run the edit-distance algorithm at the word-level instead of the character-level because most of the important variance is word-to-word. We chose the bar distinguishing high and low edit-distance by calculating the minimum edit number associated with $D$, which represents the number of insertions, deletions, and substitutions converting $C(o)$ into $C(r)$. This resulted in an edit number of 2. The edit number was chosen so that experiments using the bins of high edit-distance and low edit-distance would be sufficiently powered while ensuring the low edit-distance bucket is as close to the original prompt as possible. 

In the low-edit bin ($D \leq 2)$, we expected most of the changes to come from targeted vocabulary swaps. In order to analyze the specific words introduced by the trained attackers which resulted in over-refusal, we contrasted refused rewrites against the same attacker's complied rewrites. This effectively allows us to isolate the words which most impact the target model refusal. We used weighted log-odds with an informed Dirichlet prior following [cite] in order to rank the important added words for each model by the evidence for their contribution to over-refusal, not just their frequency.  [add the tables here properly as well as discussion of fix issues like a word appearing in all four rewrites of one prompt by requiring 3 or more distinct originals and resample originals with replacement and why we only care about the ones causally introduced by the attackers edits. also add an example? ]
\begin{figure}[tbp]
  \centering
  \includegraphics[width=\linewidth]{figures/figure2}
  \caption{\textbf{Only some over-represented words were actually introduced by the edit} Refused rewrites contrasted against the same attacker's NOT-refused rewrites in the same edit bin, so the contrast isolates refusal rather than the attacker's style. Each original counted once; words in $\geq 3$ distinct originals.}
  \label{fig:triggers}
\end{figure}[add the tables here properly as well as discussion of fix issues like a word appearing in251
all four rewrites of one prompt by requiring 3 or more distinct originals and resample originals with252
replacement and why we only care about the ones causally introduced by the attackers edits. also add253
an example? ]

For the analysis of the high edit-distance bin, we did the following. 
First, we measured how much the $\Delta(o, r)$ is displaced relative to the refusal vector from [cite: Arditi]. Then, we construct a set of 8 directions to isolate the effects of the key words we found in the previous section on over-refusal. Finally, we conduct a causal study to see how we can engineer over-refusal in models using simple word replacements/additions of the key words we found. 

[Add specifics on how exactly we fit a set of 8 directions on half of a dataset and how they are defined and what intuititvely d1 etc means]

We performed a directoinal abliteration [following arditi i think] at every layer by taking directions and replacing the residual stream with $h \rightarrow h - \langle h, d_x \rangle d_x$. We then measured over-refusal against a set of 400 held-out over-refusals and 200 genuinely harmful prompts from AdvBench [cite] in order to see if these abliterations would meaningfully change over-refusals and make models less safe on harmful prompts. 

[Show and explain table]
An interesting result of this table is that the abliteration of the key word topics reduces Llama's OR significantly, while not reducing its ability to refuse harmful prompts. This effect is most visible with the weaponization topic. 



\section{Discussion and limitations}
% Interpret what the comparison says about the refusal boundary, especially the
% gap between causal activation mechanisms and signals that are accessible
% through natural-language edits. Discuss why trained attackers and token-level
% search may find different failure modes and what model-specific trigger words
% imply for transfer.



% Cover uncertainty from independent judging, semantic-similarity models and
% safety filters; the limited set of target and attacker models; sensitivity to
% refusal openers and prompt distributions; incomplete minimal-edit evidence;
% and the possibility that some apparently benign prompts remain controversial.% Discuss dual-use risk, responsible release, and how over-refusal attacks can be
% used to improve safety--helpfulness calibration.

\section{Related work}

\subsection{Over-refusal benchmarks and prompt generation}
    \citet{rottger2024xstest} find that false refusals often relate to overfitting to superficial language features. They introduce a dataset of 250 safe prompts with homonyms and figurative language designed to trigger refusal as well as 200 legitimately unsafe contrasts. \citet{an2024phtest} utilize AutoDAN, an automatic, gradient-guided search algorithm, to create the PHTest dataset of over-refusal prompts. They find that over-refusal triggers vary by model and that increased safeguards often create increased over-refusal. \citet{cui2024orbench} scale over-refusal with ORBench, 80{,}000 seemingly toxic but benign prompts generated by small rewrites of harmful problems until they pass a judge checking for safety. \citet{wu2025evorefuse} introduce EVOREFUSE, treating prompt construction as a black box evolutionary search. They find models often over-fixate on malicious-sounding keywords while missing the surrounding context. \citet{yuan2025beyond} provide single-turn and multi-turn exaggerated-safety benchmarks, XSB. They utilize words that can be interpreted in both harmful and helpful ways to test the refusal boundary.


\subsection{Refusal representations and interventions}
    \citet{zou2023representation} identify activations that represent high-level concepts such as honesty, harmlessness, and power-seeking. They demonstrate both read and write capabilities along these vectors. \citet{marks2024geometry} investigate whether concepts such and truth and falsehood can be reduced to single vectors in activation space. They utilize difference-of-means to find these activation vectors and compare them against trained probes. \citet{arditi2024refusal} find a single direction in the residual stream that represents refusal. They show this vector is necessary for refusal by showing that refusal is deactivated by abliterating the refusal vector, projecting the residual stream onto the subspace orthogonal to the vector. They show it is sufficient by steering along the refusal vector and showing this induces refusal.
    \citet{wollschlager2025geometry} challenge the idea that every refusal is mediated by the same single axis. Using gradients to search activation space, they find a cone containing several directions that can each induce refusal. Their representational-independence test asks whether one direction still causes refusal after the effect of another has been removed; directions that survive this test behave as functionally distinct mechanisms rather than redundant descriptions of the same signal. They also test whether input prompts can actually reach these directions. This matters for us because a direction can be powerful when inserted directly into the model yet difficult to activate through small, fluent edits to the input.


\subsection{Automated red-teaming and adversarial prompt optimization}
  \citet{perez2022redteaming} use an LLM in the attacker role to automate the process of creating red-team prompts. They use an classifier on the target's response as the signal and test zero-shot generation, curated prompting, and reinforcement learning. They find reinforcement learning is most consistent but discovers a less broad array of failures. \citet{wichers2024gradient} use gradient-based methods, backpropagating from a frozen target model and safety classifier to direct a prompt rewriting algorithm. They also add a loss term based on language realism to avoid unnatural rewrites. They find rewarding fluency to be an effective improvement for creating useful red-team prompts. \citet{zou2023universal} introduce Greedy Coordinate Gradient (GCG) which iteratively optimizes an adversarial suffix based on using gradients to identify a shortlist of edits that are most likely to decrease refusal. \citet{wen2023hard} provide PEZ, a method for optimization in continuous embedding space that allows one to effectively project back onto real tokens without losing the embedding score. This is done largely through intermediate projections that prevent the embedding from straying too far from the subspace that maps onto realistic hard prompts.
  
\section{Conclusion}

TODO
% Summarize the empirical finding and its implication for secure-by-design model
% behavior without overstating small-sample or provisional results.


\appendix\section{Refusal-vector candidate selection} \label{app:refusal-vector-selection} \paragraph{Candidate construction.} For each transformer layer \(L\), we form a candidate direction \(d_L=\mu_{\mathrm{refuse},L}-\mu_{\mathrm{nonrefuse},L}\) from the activation at the final real prompt token. The contrast uses 2{,}500 prompts from each class. The prompt lists were originally divided according to Llama-3-8B-Instruct's responses; for Qwen3-32B we reuse those lists but extract Qwen activations and perform the layer selection and behavioral checks on Qwen. Each candidate is therefore model-specific even though the source prompts are shared. We use the raw difference of means because it performed better than covariance-corrected LDA in both model comparisons. \paragraph{Selection principle.} A layer can be useful in two different senses: its score may predict which prompts will be refused, or changing activations along its direction may directly change refusal. Because we use the single-layer vector as a mechanistic intervention, we prioritize causal addition and ablation evidence while using behavioral prediction as a supporting check, following the intervention-based validation of \citet{arditi2024refusal}. We do not choose a layer solely because its direction has a large norm or because it separates the same data used to construct it. \paragraph{Llama-3-8B-Instruct.} We select layer 17. Removing the layer-17 direction reduced harmful-prompt refusal from 99\% to 83\%, while adding it to benign prompts raised refusal from 0\% to 99\% at the selected coefficient. Layer 31 ranked rewrite-induced changes slightly better on the 6{,}000-pair behavioral set, but it was not causally validated; the previously considered layer 32 was nearly inert under ablation. We therefore retain layer 17 for the single causal vector and allow the multi-layer probe to use later layers when prediction, rather than intervention, is the goal. \paragraph{Qwen3-32B.} Layer 58 was the best raw-direction layer for ranking the measured change in refusal across 1{,}500 original--rewrite pairs (Spearman \(\rho=0.809\)); the multi-layer non-negative fit placed all of its weight on this layer, and the resulting score had AUC 0.976. Neighboring layers 57--61 formed a broad near-tie, so we do not interpret layer 58 as uniquely special. Adding the layer-58 direction to benign prompts raised refusal from 5\% to 93\%. The ablation test was not informative because Qwen refused only about 2\% of the harmful prompts before intervention, leaving too little baseline refusal to remove. We therefore treat layer 58 as supported by strong behavioral prediction and causal sufficiency, but not yet by causal necessity. \paragraph{Resulting choice.} These criteria yield layer 17 for Llama and layer 58 for Qwen. The asymmetric evidence should remain explicit: Llama has both ablation and addition evidence, whereas Qwen's current causal evidence is addition-only.\section{Additional experimental details}

For  generated by prompting Claude Haiku 4.5 and Claude Sonnet 4.5 to create over-refusal rewrites of benign Dolly and Alpaca prompts. Four example over-refusal rewrites are given to the Claude models as guidance. We ubegin with 5,000 benign prompts and create 3 over-refusal rewrites each to create 15,000


We sampled about 5,000 benign Dolly and Alpaca prompts. Claude Haiku 4.5 produced three meaning-preserving rewrites of each prompt using four example rewrites as guidance, giving 14,996 rewrites. Llama-3-8B-Instruct answered every rewrite and each original. We also included two smaller evaluations where either ordinary Llama or a trained Llama rewriter generated four rewrites for each of 200 held-out Alpaca prompts; plain Llama then answered them. One trained rewriter targeted the refusal direction and the other targeted our refusal probe.

Qwen (19,952 responses): One source was 1,500 original–rewrite pairs selected from a larger pool of benign Alpaca prompts rewritten by Sonnet. Qwen answered each original and rewrite four times. We also included earlier Qwen experiments where ordinary Qwen and Qwen rewriters trained using the vector, probe, or refusal-prefix signal generated rewrites of held-out Alpaca prompts; plain Qwen then answered those rewrites.

% Add acknowledgments only after double-blind review.
\bibliographystyle{plainnat}
\bibliography{references}

\end{document}


  \item \textbf{PHTest.} 